Overall, our results match part of what we would expect from the idea behind xRFM. xRFM is built for tabular data: it uses a tree to split the dataset into simpler local regions, and then learns flexible, nonlinear patterns within each region instead of relying on a fixed feature map. After preprocessing, both continuous variables and one-hot encoded categorical variables become numeric inputs, so in principle xRFM can handle a wide mix of feature types and tasks. This makes it reasonable that xRFM is often competitive with strong baselines, especially when the target depends on nonlinear interactions and there is enough data for the AGOP-based split directions to be learned reliably. At the same time, our experiments do not show xRFM as a clear overall winner. TabPFN is the most accurate model across our suite, Random Forest is usually the weakest, and xRFM and XGBoost trade places depending on the dataset. In short, xRFM performs best when it can make use of meaningful nonlinear structure in the data, but it can fall short in harder settings where its splitting and feature-learning pipeline is less stable, so its advantage over XGBoost appears real on some datasets but not consistent enough to be called a universal improvement.

The feature-importance comparison mostly supports the interpretation of AGOP given in the Methodology. On Abalone, the AGOP diagonal identifies the same broad group of important features as the other methods: weight-related variables such as \texttt{whole\_weight}, \texttt{shucked\_weight}, and \texttt{shell\_weight} are consistently ranked near the top, while the one-hot \texttt{sex\_*} features remain near the bottom. This agrees with the idea that the AGOP diagonal measures how sensitive the trained xRFM model is to each input coordinate. It is also useful that AGOP is especially close to permutation importance for the top two features, since permutation importance directly measures how much the model relies on a feature for prediction. The main differences appear in the middle of the ranking: for example, mutual information gives very high rank to \texttt{shell\_weight} and \texttt{diameter}, while PCA ranks some geometric features higher because it is driven by input variance rather than prediction. This was not a major contradiction, but it was a useful reminder that each method defines ``importance'' differently. Overall, the comparison suggests that AGOP gives a meaningful model-specific explanation of xRFM, while the disagreements with PCA and MI show that it is not merely recovering generic variance or marginal dependence patterns.

In our experiments, xRFM underperforms XGBoost most clearly on Sleep Health and Superconductor. For Sleep Health, XGBoost achieves 0.9465 Accuracy and 0.9909 AUC-ROC, compared to 0.8905 Accuracy and 0.9756 AUC-ROC for xRFM; for Superconductor, XGBoost attains a lower RMSE (9.4001 vs. 9.8139). A possible hypothesis according to xRFM’s design is that its AGOP-driven splitting becomes less reliable as the effective feature dimension grows after preprocessing, since each leaf relies on estimating a $d\times d$ AGOP matrix and extracting a leading direction from it. When $d$ is large---especially with sparse one-hot expansions (as in Sleep Health) or strong multicollinearity and redundancy among continuous descriptors (as in Superconductor)---the local feature geometry may be harder to estimate accurately, causing split directions to be noisier and downstream leaf RFMs to generalize less well. This hypothesis is consistent with the fact that the two datasets where xRFM loses most are also among the highest-dimensional settings in our suite, whereas XGBoost can more easily ignore weak or redundant dimensions through axis-aligned tree splits. That said, this evidence is correlational: confirming the mechanism would require controlled studies that vary $d$ and sparsity (e.g., ablations on one-hot cardinality or feature selection) to test whether performance degrades systematically for AGOP-based splitting in high-dimensional tabular regimes.

If we had more time, we would strengthen the empirical conclusions by broadening both the experimental scope and the analysis tools. On the scalability and performance side, we would add more datasets whose characteristics stress the regimes where xRFM may be less robust---such as higher effective dimensionality, strongly correlated continuous features, and sparse one-hot categorical expansions---to test whether the observed weaknesses are systematic rather than dataset-specific. We would also allocate a larger tuning budget to xRFM, whose computational cost limited the number of configurations explored compared to XGBoost, and search more extensively over kernel choices and bandwidth, regularization, leaf size, the use of full versus diagonal AGOP, and the number of RFM iterations; in addition, we would compare different tuning objectives (e.g., F1 and AUC-ROC instead of accuracy for potentially imbalanced classification tasks, and alternative regression losses beyond RMSE) to better match evaluation goals. Finally, for interpretability, we would extend the comparison beyond PCA/MI/permutation by including methods such as SHAP and by contrasting against feature importance derived from strong tree baselines (Random Forest/XGBoost), helping to clarify whether AGOP provides a genuinely distinct, model-specific perspective or largely agrees with established explanations.
